% A transformação homogênea entre os referenciais consecutivos $\{i-1\}$ e $\{i\}$ é definida pela convenção de \gls{dh} \cite{craig2005}.
% Cada transformação descreve a relação geométrica entre dois referenciais adjacentes e depende dos quatro parâmetros apresentados na Tabela~\ref{tab_dh_zxx}.

% A matriz de transformação (Equação~\ref{eq_Ti_1}) é construída pela composição de quatro operações fundamentais:
% uma rotação de $\alpha_{i-1}$ em torno de $\hat{X}_{i-1}$, seguida de uma translação de $a_{i-1}$ ao longo de $\hat{X}_{i-1}$, uma rotação de $\theta_i$ em torno de $\hat{Z}_i$ e, por fim, uma translação de $d_i$ ao longo de $\hat{Z}_i$ \cite{craig2005}.
% Assim, utilizando as matrizes elementares de rotação e translação, tem-se:

% \begin{equation}
% {}^{i-1}T_i
% =
% T_X(\alpha_{i-1})\,D_X(a_{i-1})\,T_Z(\theta_i)\,D_Z(d_i).
% \label{eq_Ti_1}
% \end{equation}

% Portanto, a partir da Equação~\ref{eq_Ti_1}, resulta em:

% \begin{equation}
% {}^{i-1}_{\ \ \ i}T(q_i)
% =
% {}^{i-1}_{\ \ \ i'}A \, {}^{i'}_{\ \ i}A
% =
% \begin{bmatrix}
% \cos\theta_i & -\sin\theta_i \cos\alpha_i & \sin\theta_i \sin\alpha_i & a_i \cos\theta_i \\
% \sin\theta_i & \cos\theta_i \cos\alpha_i & -\cos\theta_i \sin\alpha_i & a_i \sin\theta_i \\
% 0 & \sin\alpha_i & \cos\alpha_i & d_i \\
% 0 & 0 & 0 & 1
% \end{bmatrix}.
% \end{equation}

% O termo $T_X(\alpha_{i-1})$ representa a rotação de ângulo $\alpha_{i-1}$ em torno de $\hat{X}_{i-1}$,
% $D_X(a_{i-1})$ representa a translação de magnitude $a_{i-1}$ ao longo de $\hat{X}_{i-1}$,
% $T_Z(\theta_i)$ representa a rotação de ângulo $\theta_i$ em torno de $\hat{Z}_i$ e
% $D_Z(d_i)$ representa a translação de magnitude $d_i$ ao longo de $\hat{Z}_i$.
% %%%%%%%%%%%%%%%%%%%%%%%%%%%%%%%%%%%
% As quatro matrizes elementares utilizadas na convenção de \gls{dh} podem ser escritas como \cite{craig2005}:

% \begin{equation}
% \label{eq_tx}
% T_X(\alpha_{i-1}) =
% \begin{bmatrix}
% 1 & 0 & 0 & 0 \\
% 0 & \cos\alpha_{i-1} & -\sin\alpha_{i-1} & 0 \\
% 0 & \sin\alpha_{i-1} &  \cos\alpha_{i-1} & 0 \\
% 0 & 0 & 0 & 1
% \end{bmatrix},
% \end{equation}

% \begin{equation}
% D_X(a_{i-1}) =
% \begin{bmatrix}
% 1 & 0 & 0 & a_{i-1} \\
% 0 & 1 & 0 & 0 \\
% 0 & 0 & 1 & 0 \\
% 0 & 0 & 0 & 1
% \end{bmatrix},
% \end{equation}

% \begin{equation}
% T_Z(\theta_i) =
% \begin{bmatrix}
% \cos\theta_i & -\sin\theta_i & 0 & 0 \\
% \sin\theta_i &  \cos\theta_i & 0 & 0 \\
% 0 & 0 & 1 & 0 \\
% 0 & 0 & 0 & 1
% \end{bmatrix},
% \end{equation}

% \begin{equation}
%     \label{eq_Dz}
% D_Z(d_i) =
% \begin{bmatrix}
% 1 & 0 & 0 & 0 \\
% 0 & 1 & 0 & 0 \\
% 0 & 0 & 1 & d_i \\
% 0 & 0 & 0 & 1
% \end{bmatrix}.
% \end{equation}



% %%%%%%%%%%%%%%%%%%%%%%%%%%%%

% A matriz de transformação homogênea pode ser particionada em duas submatrizes: a matriz de rotação ${}^{i-1}R_i$ ($3 \times 3$), que descreve a orientação do referencial $\{i\}$ em relação a $\{i-1\}$, e o vetor de posição ${}^{i-1}P_i$ ($3 \times 1$), que descreve a posição da origem de $\{i\}$ expressa em $\{i-1\}$.
% Conforme \cite{craig2005}, essa estrutura é escrita como:

% \begin{equation}
% {}^{i-1}T_i =
% \begin{bmatrix}
% {}^{i-1}R_i & {}^{i-1}P_i \\
% \mathbf{0}_{1\times 3} & 1
% \end{bmatrix}.
% \label{eq_estrutura_submatrizes}
% \end{equation}

% %%%%%%%%%%%%%%%%%%%%%%%%%%%%%%%%%%%%

% Para o \gls{mr} em estudo, composto por três juntas revolutas ($i=1,2,3$), as transformações homogêneas entre elos consecutivos são obtidas pela substituição direta dos parâmetros da Tabela~\ref{tab_dh_zxx} na forma geral, com $d_i=0$. Assim, as matrizes de transformação para cada elo são:

% %%%%%%%%%%%% 0 a 1
% \begin{equation}
% {}^{0}T_1 =
% \begin{bmatrix}
% \cos\theta_1 & 0 & \sin\theta_1 & 0 \\
% \sin\theta_1 & 0 & -\cos\theta_1 & 0 \\
% 0            & 1 & 0             & 0 \\
% 0            & 0 & 0             & 1
% \end{bmatrix},
% \label{eq_T01}
% \end{equation}

% %%%%%%%%%%%% 1 a 2
% \begin{equation}
% {}^{1}T_2 =
% \begin{bmatrix}
% \cos\theta_2 & -\sin\theta_2 & 0 & a_2\cos\theta_2 \\
% \sin\theta_2 & \cos\theta_2  & 0 & a_2\sin\theta_2 \\
% 0            & 0             & 1 & 0 \\
% 0            & 0             & 0 & 1
% \end{bmatrix},
% \label{eq_T12}
% \end{equation}

% %%%%%%%%%%%% 2 a 3
% \begin{equation}
% {}^{2}T_3 =
% \begin{bmatrix}
% \cos\theta_3 & -\sin\theta_3 & 0 & a_3\cos\theta_3 \\
% \sin\theta_3 & \cos\theta_3  & 0 & a_3\sin\theta_3 \\
% 0            & 0             & 1 & 0 \\
% 0            & 0             & 0 & 1
% \end{bmatrix},
% \label{eq_T23}
% \end{equation}


% Definindo $\theta_{23}=\theta_2+\theta_3$, a transformação homogênea acumulada do referencial $\{0\}$ ao referencial $\{3\}$ é dada por:

% \begin{equation}
% {}^{0}T_{3} = {}^{0}T_{1}\,{}^{1}T_{2}\,{}^{2}T_{3}.
% \label{eq_T03_def}
% \end{equation}

% Substituindo \eqref{eq_T01}--\eqref{eq_T23} em \eqref{eq_T03_def} e efetuando as multiplicações, obtém-se:

% \begin{equation}
% {}^{0}T_{3} =
% \begin{bmatrix}
% \cos\theta_{1}\cos\theta_{23} & -\cos\theta_{1}\sin\theta_{23} & \sin\theta_{1} & \cos\theta_{1}\bigl(a_2\cos\theta_{2}+a_3\cos\theta_{23}\bigr)\\
% \sin\theta_{1}\cos\theta_{23} & -\sin\theta_{1}\sin\theta_{23} & -\cos\theta_{1} & \sin\theta_{1}\bigl(a_2\cos\theta_{2}+a_3\cos\theta_{23}\bigr)\\
% \sin\theta_{23} & \cos\theta_{23} & 0 & a_2\sin\theta_{2}+a_3\sin\theta_{23}\\
% 0 & 0 & 0 & 1
% \end{bmatrix}.
% \label{eq_T03_final}
% \end{equation}

% A inversa de uma transformação homogênea ${}^{0}T_{3}=\begin{bmatrix}{}^{0}R_{3} & {}^{0}P_{3}\\ \mathbf{0}_{1\times 3} & 1\end{bmatrix}$ é dada por:

% \begin{equation}
% {}^{3}T_{0} = ({}^{0}T_{3})^{-1}
% =
% \begin{bmatrix}
% ({}^{0}R_{3})^{T} & -({}^{0}R_{3})^{T}\,{}^{0}P_{3}\\
% \mathbf{0}_{1\times 3} & 1
% \end{bmatrix}.
% \label{eq_T30_geral}
% \end{equation}

% Aplicando \eqref{eq_T30_geral} ao resultado de \eqref{eq_T03_final}, obtém-se:

% \begin{equation}
% {}^{3}T_{0} =
% \begin{bmatrix}
% \cos\theta_{1}\cos\theta_{23} & \sin\theta_{1}\cos\theta_{23} & \sin\theta_{23} & -a_2\cos\theta_{3}-a_3\\
% -\cos\theta_{1}\sin\theta_{23} & -\sin\theta_{1}\sin\theta_{23} & \cos\theta_{23} & -a_2\sin\theta_{3}\\
% \sin\theta_{1} & -\cos\theta_{1} & 0 & 0\\
% 0 & 0 & 0 & 1
% \end{bmatrix}.
% \label{eq_T30_final}
% \end{equation}

% %%%%%%%%%%%%%%%%%%%%%%%%%

% Essa matriz descreve a transformação homogênea do referencial $\{0\}$ em relação ao referencial $\{3\}$, obtida a partir da inversa da transformação apresentada em \eqref{eq_T03_final}.


%%%%%%%%%%%% após 20/04/2026:

Como apresentado na Figura~\ref{fig_rotacaoJuntas_MR}, o \gls{mr} segue a mesma construção adotada para o braço planar descrito por \cite{sciavicco_siciliano2010}. Com essa definição, as matrizes de transformação homogênea associadas às juntas do \gls{mr} são dadas por:

\begin{equation}
{}^{0}_{1}T
=
\begin{bmatrix}
\cos\theta_{1} & 0 & \sin\theta_{1} & 0 \\
\sin\theta_{1} & 0 & -\cos\theta_{1} & 0 \\
0 & 1 & 0 & 0 \\
0 & 0 & 0 & 1
\end{bmatrix}.
\label{eq_T01}
\end{equation}

Para $i=2$, tem-se:

\begin{equation}
{}^{1}_{2}T=
\begin{bmatrix}
\cos\theta_{2} & -\sin\theta_{2} & 0 & a_{2}\cos\theta_{2} \\
\sin\theta_{2} & \cos\theta_{2} & 0 & a_{2}\sin\theta_{2} \\
0 & 0 & 1 & 0 \\
0 & 0 & 0 & 1
\end{bmatrix}.
\label{eq_T12_antropomorfico}
\end{equation}

Por sua vez, para $i=3$, escreve-se:

\begin{equation}
{}^{2}_{3}T=
\begin{bmatrix}
\cos\theta_{3} & -\sin\theta_{3} & 0 & a_{3}\cos\theta_{3} \\
\sin\theta_{3} & \cos\theta_{3} & 0 & a_{3}\sin\theta_{3} \\
0 & 0 & 1 & 0 \\
0 & 0 & 0 & 1
\end{bmatrix}.
\label{eq_T23_antropomorfico}
\end{equation}

A cinemática direta do \gls{mr} é obtida pela composição sucessiva dessas transformações homogêneas.
Assim, para o vetor de coordenadas generalizadas:

\begin{equation}
q=
\begin{bmatrix}
\theta_{1} & \theta_{2} & \theta_{3}
\end{bmatrix}^{T}.
\end{equation}

Obtém-se: 

\begin{equation}
{}^{0}_{3}T(q)={}^{0}_{1}T\,{}^{1}_{2}T\,{}^{2}_{3}T.
\label{eq_T03_def}
\end{equation}

Definindo $\theta_{23}=\theta_{2}+\theta_{3}$, a multiplicação das matrizes em \eqref{eq_T03_def} conduz a:

\begin{equation}
{}^{0}_{3}T(q)=
\begin{bmatrix}
\cos\theta_{1}\cos\theta_{23} & -\cos\theta_{1}\sin\theta_{23} & \sin\theta_{1} & \cos\theta_{1}\left(a_{2}\cos\theta_{2}+a_{3}\cos\theta_{23}\right) \\
\sin\theta_{1}\cos\theta_{23} & -\sin\theta_{1}\sin\theta_{23} & -\cos\theta_{1} & \sin\theta_{1}\left(a_{2}\cos\theta_{2}+a_{3}\cos\theta_{23}\right) \\
\sin\theta_{23} & \cos\theta_{23} & 0 & a_{2}\sin\theta_{2}+a_{3}\sin\theta_{23} \\
0 & 0 & 0 & 1
\end{bmatrix}.
\label{eq_T03_final}
\end{equation}

A Equação~\eqref{eq_T03_final} representa a transformação homogênea entre o referencial da espaçonave e o referencial $\{3\}$, fornecendo a orientação e a posição do último elo em função das variáveis articulares.
Observa-se, ainda, que o eixo $z_{3}$ permanece alinhado com $z_{2}$ \cite{sciavicco_siciliano2010}.
